\subsection{window-\texorpdfstring{$6$}{6} 的中心 Grassmann 对称、交换化轨道与 boundary 不可恢复性}\label{subsec:conclusion-window6-center-grassmann-abelian-boundary-nonrecoverability}
在 window-\(6\) 的 fold-gauge 全自同构群、中心直和分裂、交换化 parity-charge 分解以及纤维群胚代数的 Wedderburn 分解都已固定之后，boundary 三维子格的几何规范性还可进一步压缩为一个统一的不可恢复性原理：从抽象群层面看，整个 dyadic 中心只是一个单一的 \(8\) 维 \(\FF_2\) 自然模；从交换化层面看，全部可见轨道只记忆二点/三点/四点纤维部门的维数与 Hamming 重数，而不记忆 boundary 的具体选址；从仅保留抽象半单代数结构的层面看，全部 \(M_2(\CC)\) 简单块同样只以块数出现，因而 boundary 的 \(3+5\) 分割不能由代数同构型本身恢复。

\begin{theorem}[中心 dyadic 模的特征子群刚性]\label{thm:conclusion-window6-center-characteristic-rigidity}
设
\[
G:=\mathcal G_6^{\mathrm{bin}}
\cong
(\ZZ/2\ZZ)^8\times S_3^{\,4}\times S_4^{\,9}.
\]
则任一满足 \(W\le Z(G)\) 且在 \(\Aut(G)\) 下不变的子群都只能是
\[
W=0
\qquad\text{或}\qquad
W=Z(G).
\]
特别地，由定理 \ref{thm:window6-foldbin-gauge-center-boundary-direct-sum} 给出的规范边界子群
\[
Z_{\partial}\cong (\ZZ/2\ZZ)^3
\]
不是 \(G\) 的特征子群。
\end{theorem}
\begin{proof}
由定理 \ref{thm:terminal-window6-fold-gauge-aut-out-structure}，
\[
G=A\times H,
\qquad
A:=Z(G)\cong (\ZZ/2\ZZ)^8,
\qquad
H:=S_3^{\,4}\times S_4^{\,9},
\]
并且对任意 \(\alpha\in\Aut(A)\cong \GL(8,2)\) 都存在自同构
\[
\varphi_\alpha(a,h):=(\alpha(a),h).
\]
因此任一在 \(\Aut(G)\) 下不变的子群 \(W\le Z(G)=A\) 必在 \(\GL(8,2)\) 的自然作用下不变。若 \(W\neq 0\)，取 \(0\neq v\in W\)。由于 \(\GL(8,2)\) 在 \(\FF_2^8\setminus\{0\}\) 上传递，\(W\) 必包含全部非零向量，从而 \(W=A=Z(G)\)。故除 \(0\) 与 \(Z(G)\) 外不存在其他 \(\Aut(G)\)-不变子群。

最后，定理 \ref{thm:window6-foldbin-gauge-center-boundary-direct-sum} 给出 \(Z_{\partial}\) 是 \(Z(G)\) 的一个非零真三维子群，故不可能为特征子群。证毕。
\end{proof}

\begin{corollary}[中心商的 Grassmann 塌缩]\label{cor:conclusion-window6-center-quotient-grassmann-collapse}
沿用定理 \ref{thm:conclusion-window6-center-characteristic-rigidity} 的记号。对任意 \(0\le k\le 8\) 与任意两个 \(k\)-维子空间
\[
U,V\le Z(G)\cong \FF_2^8,
\]
都存在 \(\varphi\in \Aut(G)\) 使 \(\varphi(U)=V\)。因而
\[
G/U\cong G/V\cong (\ZZ/2\ZZ)^{8-k}\times S_3^{\,4}\times S_4^{\,9}.
\]
特别地，boundary 商群 \(G/Z_{\partial}\) 与中心中任意其他三维子空间所对应的商群完全同构；boundary 几何在抽象中心商层面不可辨识。
\end{corollary}
\begin{proof}
由定理 \ref{thm:conclusion-window6-center-characteristic-rigidity} 的证明，\(\Aut(G)\) 在 \(Z(G)\) 上诱导整个 \(\GL(8,2)\) 的自然作用。于是 \(\GL(8,2)\) 在 Grassmann 流形
\[
\mathrm{Gr}(k,8;\FF_2)
\]
上作用传递，故存在某个 \(\varphi\in\Aut(G)\) 满足 \(\varphi(U)=V\)。这立即推出 \(G/U\cong G/V\)。

另一方面，仍由定理 \ref{thm:terminal-window6-fold-gauge-aut-out-structure}，
\[
G=Z(G)\times H,
\qquad
H:=S_3^{\,4}\times S_4^{\,9},
\]
故
\[
G/U\cong Z(G)/U\times H\cong (\ZZ/2\ZZ)^{8-k}\times S_3^{\,4}\times S_4^{\,9}.
\]
证毕。
\end{proof}

\begin{theorem}[parity-charge 交换化的自同构轨道分类]\label{thm:conclusion-window6-abelianization-aut-orbits}
设
\[
G:=\mathcal G_6^{\mathrm{bin}},
\]
并在定理 \ref{thm:window6-foldbin-gauge-center-boundary-direct-sum} 的分解下记
\[
A:=Z_{\partial}\oplus Z_{\mathrm{int}}\cong \FF_2^8,
\qquad
B_3:=\bigl(S_3^{\mathrm{ab}}\bigr)^4\cong \FF_2^4,
\qquad
B_4:=\bigl(S_4^{\mathrm{ab}}\bigr)^9\cong \FF_2^9.
\]
则
\[
G^{\mathrm{ab}}\cong A\oplus B_3\oplus B_4.
\]
\(\Aut(G)\) 在 \(G^{\mathrm{ab}}\) 上的轨道恰由下列数据完全分类：
\begin{enumerate}
  \item 零元轨道 \(\{0\}\)；
  \item 非零二点纤维中心轨道 \(A\setminus\{0\}\)；
  \item 对每个满足 \(0\le i\le 4\)、\(0\le j\le 9\)、\((i,j)\neq(0,0)\) 的整数对，恰有一个轨道
  \[
  \mathcal O_{i,j}
  :=
  \left\{
  (a,u,v)\in A\oplus B_3\oplus B_4:
  \operatorname{wt}(u)=i,\ \operatorname{wt}(v)=j,\ (u,v)\neq(0,0)
  \right\}.
  \]
\end{enumerate}
因此轨道总数恰为
\[
2+(5\cdot 10-1)=51.
\]
\end{theorem}
\begin{proof}
内自同构在交换化上作用平凡，故只需分析定理 \ref{thm:xi-window6-full-automorphism-group-structure} 给出的外部作用。该定理把 \(\Aut(G)\) 写成
\[
\Hom\!\bigl((\ZZ/2\ZZ)^{13},(\ZZ/2\ZZ)^8\bigr)\rtimes
\Bigl(
\GL_8(\FF_2)\times (S_3\wr S_4)\times (S_4\wr S_9)
\Bigr).
\]
传到交换化后，两个 wreath product 的逐块内自同构部分都消失，只剩
\[
S_4\times S_9
\]
分别在 \(B_3\cong \FF_2^4\) 与 \(B_4\cong \FF_2^9\) 的标准坐标上作全排列。因此任意自同构在
\[
G^{\mathrm{ab}}=A\oplus B_3\oplus B_4
\]
上都可写成
\[
(a,u,v)\longmapsto \bigl(\alpha a+T(u,v),\ \sigma u,\ \tau v\bigr),
\]
其中
\[
\alpha\in \GL(A),\qquad
T\in \Hom(B_3\oplus B_4,A),\qquad
\sigma\in S_4,\qquad
\tau\in S_9.
\]

若 \((u,v)=(0,0)\)，则上式退化为 \(A\) 上的自然 \(\GL(8,2)\)-作用；由其在 \(A\setminus\{0\}\) 上传递，得到零元轨道与唯一的非零二点纤维中心轨道。

若 \((u,v)\neq(0,0)\)，则可取线性泛函 \(\lambda:B_3\oplus B_4\to \FF_2\) 使 \(\lambda(u,v)=1\)，并定义
\[
T(x):=\lambda(x)\,a\qquad (x\in B_3\oplus B_4).
\]
于是 \(T(u,v)=a\)，从而 \((a,u,v)\) 与 \((0,u,v)\) 属同一轨道。接着 \(S_4\) 与 \(S_9\) 分别在 \(\FF_2^4\) 与 \(\FF_2^9\) 的固定 Hamming 重层上作用传递，故轨道只由 \(\operatorname{wt}(u)\) 与 \(\operatorname{wt}(v)\) 决定；反向方向则因坐标排列不改变 Hamming 重而显然。遂得所述全部轨道，并得到总数
\[
2+(5\cdot 10-1)=51.
\]
证毕。
\end{proof}

\begin{theorem}[抽象半单代数口径下的 boundary 遗忘]\label{thm:conclusion-window6-semisimple-boundary-forgetfulness}
设
\[
\mathcal A_6:=\CC[\mathcal R_6^{\mathrm{bin}}]
\cong
M_2(\CC)^{\oplus 8}\oplus M_3(\CC)^{\oplus 4}\oplus M_4(\CC)^{\oplus 9}.
\]
令 \(I_{\partial}\subset \mathcal A_6\) 为由三条 boundary 二点纤维所对应简单块张成的双边理想，则
\[
I_{\partial}\cong M_2(\CC)^{\oplus 3}.
\]
并且 \(\mathrm{Aut}_{\CC\text{-alg}}(\mathcal A_6)\) 在全部与 \(M_2(\CC)^{\oplus 3}\) 同构的双边理想集合上作用传递；该轨道大小恰为
\[
\binom{8}{3}=56.
\]
因此，若只保留 \(\mathcal A_6\) 的抽象复半单代数结构，则 boundary 三块与任意其他三块 dyadic \(M_2\)-因子在代数上不可区分；特别地，dyadic 部门内部的 \(3+5\) boundary/internal 分割不能由 \(\mathcal A_6\) 的同构型本身恢复。
\end{theorem}
\begin{proof}
Wedderburn 分解与 window-\(6\) 直方图给出
\[
\mathcal A_6
\cong
M_2(\CC)^{\oplus 8}\oplus M_3(\CC)^{\oplus 4}\oplus M_4(\CC)^{\oplus 9},
\]
而三条 boundary 纤维都对应 dyadic 简单块，故 \(I_{\partial}\cong M_2(\CC)^{\oplus 3}\)。

再由定理 \ref{thm:conclusion-foldbin-groupoid-aut-pgl-pu} 在 \(m=6\) 的专门化，
\[
\mathrm{Aut}_{\CC\text{-alg}}(\mathcal A_6)
\cong
\bigl(\operatorname{PGL}_2(\CC)^8\rtimes S_8\bigr)
\times
\bigl(\operatorname{PGL}_3(\CC)^4\rtimes S_4\bigr)
\times
\bigl(\operatorname{PGL}_4(\CC)^9\rtimes S_9\bigr).
\]
其中 \(S_8\) 在八个 \(M_2(\CC)\)-块上作全置换。故对任意两个三元子集 \(S,T\subset\{1,\dots,8\}\)，若 \(|S|=|T|=3\)，总可取 \(\pi\in S_8\) 使 \(\pi(S)=T\)；对应自同构便把
\[
\bigoplus_{i\in S} M_2(\CC)
\]
送到
\[
\bigoplus_{i\in T} M_2(\CC).
\]
因此全部 \(M_2(\CC)^{\oplus 3}\)-型双边理想构成单一轨道，轨道大小即三元子集总数 \(\binom{8}{3}=56\)。特别地，boundary 所对应的三块在抽象半单代数层面并不内在。证毕。
\end{proof}

\begin{corollary}[交换子复杂度的非 dyadic 局域化]\label{cor:conclusion-window6-commutator-support-nondyadic}
设
\[
G:=\mathcal G_6^{\mathrm{bin}}
\cong
(\ZZ/2\ZZ)^8\times S_3^{\,4}\times S_4^{\,9}.
\]
则其导出列满足
\[
G'
\cong
C_3^{\,4}\times A_4^{\,9},
\qquad
G''
\cong
V_4^{\,9},
\qquad
G^{(3)}=1.
\]
因此 \(G\) 的导出长度恰为 \(3\)。特别地，全部非平凡交换子复杂度都集中在三点与四点纤维部门，而二点纤维部门只以中心电荷
\[
Z(G)\cong (\ZZ/2\ZZ)^8
\]
的形式存留。
\end{corollary}
\begin{proof}
定理 \ref{thm:conclusion-foldbin-gauge-derived-fitting-composition} 已给出
\[
G'
\cong
A_3^{\,4}\times A_4^{\,9},
\qquad
G''
\cong
V_4^{\,9},
\qquad
G^{(3)}=1.
\]
由于 \(A_3\cong C_3\)，得第一式。另一方面，\(G\) 的 dyadic 因子 \((\ZZ/2\ZZ)^8\cong S_2^{\,8}\) 本身阿贝尔，并由定理 \ref{thm:xi-foldbin-center-solvable-fiber-criterion} 正好构成全部中心；故它不对任一高阶交换子群作出非平凡贡献。证毕。
\end{proof}

\noindent
由此，window-\(6\) 的 boundary 结构在几何建模时仍是规范的三轴对象；但一旦投向抽象群、自同构轨道或仅保留同构型的半单代数，同一对象便退化为纯粹的维数型信息：中心只记忆 \(k\)-平面的维数，交换化只记忆二点/三点/四点部门的重数与 Hamming 重对，而群胚代数只记忆块大小直方图 \(2:8,\ 3:4,\ 4:9\)。因此，boundary/internal 的具体选址必须依赖额外的坐标、记录轴或 \(K_0\) 支撑数据，而不能由这些抽象对象单独反演。

\endinput
